\documentclass{article}
\usepackage[UTF8]{ctex}
\usepackage{amsmath,amssymb,amsfonts}
\usepackage{enumitem}
\usepackage{geometry}
\geometry{a4paper, margin=1in}
% === 错题本排版优化（不参与题目标识） ===
\setlength{\parindent}{0pt}
\setlength{\parskip}{0.35em plus 0.10em minus 0.08em}
\setlength{\emergencystretch}{3em}
\setlist{itemsep=0.25em, topsep=0.40em, parsep=0pt}
\allowdisplaybreaks[2]
% === 错题本排版优化结束 ===
\begin{document}
\section*{南京航空航天大学2026年数学分析真题}
\section*{100．南京航空航天大学2026年数学分析试卷（回忆版）}

\subsection{题目1}
1．计算题．

\begin{itemize}
  \item[（1）] 求极限 \(\lim _{x \rightarrow \infty}\left[\left(x^{3}-\frac{1}{2} x^{2}-x\right) e^{\frac{1}{x}}-\sqrt{x^{6}+1}\right]\) 。
  \item[（2）] 求积分 \(\int_{-\frac{\pi}{4}}^{\frac{\pi}{4}} \frac{\tan ^{2} x}{1+e^{-x}} \mathrm{~d} x\) ．
\end{itemize}

\subsection{题目2}
2．设 \(f(x)\) 在 \([0,1]\) 上可微，\(f(1)-2 \int_{0}^{\frac{1}{2}} x f(x) \mathrm{d} x=0\) ，证明：存在 \(\xi \in(0,1)\) ，使得 \(f^{\prime}(\xi)=-\frac{f(\xi)}{\xi}\) ．

\subsection{题目3}
3．压缩数列 \(\left|x_{n}-x_{n-1}\right| \leq r\left|x_{n-1}-x_{n-2}\right|, r \in(0,1)\) ，证明：\(\left\{x_{n}\right\}\) 一定收敛．

\subsection{题目4}
4．解答如下问题：

\begin{itemize}
  \item[（1）] 未知。
  \item[（2）] \(f(x)=\frac{x+2}{x+1}, a_{1}=1, a_{n+1}=f\left(a_{n}\right)\) ，证明：\(\left\{a_{n}\right\}\) 收敛，并求极限．
\end{itemize}

\subsection{题目5}
5．证明导数有界则一致连续，反之，一致连续导数是否有界？证明或给出反例．

\subsection{题目6}
6．解答如下问题：

\begin{itemize}
  \item[（1）] 判断 \(\int_{1}^{+\infty} \frac{\sin x}{x^{\alpha}+\sin x} \mathrm{~d} x\) 的敛散性．
  \item[（2）] 证明：\(f(x)=x e^{-x^{2}} \int_{0}^{x} e^{-t^{2}} \mathrm{~d} t\) 在 \([0,+\infty)\) 上一致连续．
\end{itemize}

\subsection{题目7}
7．设 \(f(x), g(x)\) 在 \([a, b]\) 上连续．

\begin{itemize}
  \item[（1）] \([f(x)-f(y)][g(x)-g(y)] \geq 0\) ，证明： \(\int_{a}^{b} f(x) \mathrm{d} x \int_{a}^{b} g(x) \mathrm{d} x \geq(b-a) \int_{a}^{b} f(x) g(x) \mathrm{d} x\) ．
  \item[（2）] \(\int_{0}^{1} \frac{f(x)}{1-f(x)} \leq \frac{\int_{0}^{1} f(x) \mathrm{d} x}{\int_{0}^{1}[1-f(x)] \mathrm{d} x}\) ．
\end{itemize}

\subsection{题目8}
8．证明 \(f(x)=\sum_{n=1}^{\infty} \frac{\left|x-\frac{1}{n}\right|}{2^{n}}\) 在 \(x_{k}=\frac{1}{k}(k=1,2, \cdots)\) 点不可微，求导函数．

\subsection{题目9}
9．已知

\[
f(x, y)= \begin{cases}x-y+\frac{x y^{2}}{x^{2}+y^{2}}, & x^{2}+y^{2} \neq 0 \\ 0, & x^{2}+y^{2}=0\end{cases}
\]

求 \(f(x, y)\) 偏导存在可微性连续性 \((? ? ?), f(x, y)\) 在 \((0,0)\) 点沿方向 \(\mathbf{l}=(\cos \theta, \sin \theta)\) 的方向导数．

\subsection{题目10}
10．\(z=z(x, y)\) 存在二阶偏导，\(u=x^{2}-y^{2}, v=2 x y, z_{x x}+z_{y y}=0\) ，用 \(u, v\) 自变量表示。

\subsection{题目11}
11．有一球体半径为 \(a\) ，球外有一点 \(M(0,0, b)(b>a)\) ，从点 \(M\) 作与球相切的圆锥，求球与圆雉所围立体的体积。

\subsection{题目12}
12．求曲面积分

\[
\iint_{\sum} y^{2} z \mathrm{~d} x \mathrm{~d} y+x z \mathrm{~d} x \mathrm{~d} z+x y^{2} \mathrm{~d} y \mathrm{~d} z .
\]

\(\sum\) 为 \(x^{2}+y^{2}=z, x^{2}+y^{2}=1\) 与坐标平面围成立体的表面，取外侧．
\end{document}
